
\section{Introduction to SQL}

\begin{frame}{\textcolor{brown}{SQL \& CRUD}}
  \begin{block}{SQL — Structured Query Language}
    The \textbf{standard language} for communicating with relational databases.
    SQL lets you perform four fundamental operations, collectively called \textbf{CRUD}:
  \end{block}

  \vspace{0.5em}
  \begin{center}
    \renewcommand{\arraystretch}{1.3}
    \begin{tabular}{ccll}
      \toprule
      \textbf{Letter} & \textbf{Operation} & \textbf{SQL Command} & \textbf{Purpose} \\
      \midrule
      \textbf{C} & Create & \texttt{INSERT}  & Add new records \\
      \textbf{R} & Read   & \texttt{SELECT}  & Retrieve records \\
      \textbf{U} & Update & \texttt{UPDATE}  & Modify existing records \\
      \textbf{D} & Delete & \texttt{DELETE}  & Remove records \\
      \bottomrule
    \end{tabular}
  \end{center}

  \vspace{0.5em}
  % \begin{alertblock}{\faExclamationCircle\; Syntax Rule}
\emph{\textbf{NB: }}All SQL statements must end with a semicolon \textbf{`` \texttt{;} ''}
  % \end{alertblock}
\end{frame}

% ─────────────────────────────────────────────────────────────────────────────
\begin{frame}[fragile]{\textcolor{brown}{Creating a Database \& Your First Table}}
Create a database:
\lstinputlisting[style=sql]{../src/component/src/create_db.sql}
% \begin{lstlisting}[style=sql]
% CREATE DATABASE school;
% \end{lstlisting}

\begin{block}{Create a table to store student records:}
\lstinputlisting[style=sql]{../src/component/src/create_table.sql}
\end{block}
  \begin{columns}[T]
    \column{0.46\textwidth}
      \textbf{Data Types}
      \begin{itemize}
        \item \texttt{SERIAL} — auto-incrementing integer
        \item \texttt{VARCHAR(n)} — string up to \textit{n} characters
        \item \texttt{INTEGER} — whole number
        \item \texttt{BOOLEAN} — true / false
      \end{itemize}

    \column{0.50\textwidth}
      \textbf{Constraints}
      \begin{itemize}
        \item \texttt{PRIMARY KEY} — uniquely identifies each row
        \item \texttt{NOT NULL} — field is required; cannot be empty
        \item \texttt{UNIQUE} — no duplicate values allowed
        \item \texttt{DEFAULT} — sets a fallback value
      \end{itemize}
  \end{columns}
\end{frame}